% ============================================================================
%  FORM-TL-05 - Taller de Revision Tecnica Cruzada del Codigo del PFC
%  Equipo revisor: FUVV  |  PFC revisado: ACC - Soporte Tecnico ISP
%  Compilacion: pdflatex informe_revision.tex  (x3, sin biber)
% ============================================================================
\documentclass[10pt,a4paper]{article}

\usepackage[utf8]{inputenc}
\usepackage[T1]{fontenc}
% lmodern omitido: no es necesario y evita dependencias extra
\usepackage[margin=2.1cm]{geometry}
\usepackage{amssymb}
\usepackage{booktabs}
\usepackage{tabularx}
\usepackage{xltabular}
\usepackage{longtable}
\usepackage{array}
\usepackage{caption}
\usepackage{enumitem}
\usepackage{listings}
\usepackage{xcolor}
\usepackage{tcolorbox}
\usepackage{xurl}
\usepackage[hidelinks]{hyperref}

% --- Nomenclatura en espanol sin babel (garantiza compilacion en clon limpio) --
\renewcommand{\figurename}{Figura}
\renewcommand{\tablename}{Tabla}
\renewcommand{\refname}{Referencias}
\renewcommand{\contentsname}{Indice}
\renewcommand{\lstlistingname}{Fragmento}

% --------------------------- DATOS EDITABLES --------------------------------
\newcommand{\repoRevisado}{https://github.com/carlospatroner-boop/PFC-SOPORTE-ISP}
\newcommand{\repoRevisor}{https://github.com/PENDIENTE/FUVV-revision-tl05}
\newcommand{\commitCorto}{479b637}
\newcommand{\commitLargo}{479b6370a640c0d2edfad801e6937e9ede8c3c80}
\newcommand{\fechaAuditoria}{Martes 11 de agosto de 2026, 21h30 (UTC-5)}
\newcommand{\cierreSemana}{Viernes 14 de agosto de 2026}
% ----------------------------------------------------------------------------

\definecolor{codebg}{RGB}{247,247,247}
\definecolor{coderule}{RGB}{210,210,210}
\lstset{
  basicstyle=\ttfamily\scriptsize,
  backgroundcolor=\color{codebg},
  frame=single,
  rulecolor=\color{coderule},
  breaklines=true,
  columns=fullflexible,
  numbers=left,
  numberstyle=\ttfamily\tiny\color{gray},
  xleftmargin=2.2em,
  framexleftmargin=2.0em,
  captionpos=b,
  showstringspaces=false
}

\newcommand{\sev}[1]{\textbf{#1}}
\DeclareUrlCommand\ruta{\urlstyle{tt}}
\newcount\cntlineas
\newcommand{\lineas}[1]{\leavevmode\hfill\par\vspace{3pt}\cntlineas=#1\relax
  \loop\ifnum\cntlineas>0
    \noindent\rule{\linewidth}{0.4pt}\par\vspace{9pt}%
    \advance\cntlineas by -1 \repeat}

\renewcommand{\arraystretch}{1.15}
\setlength{\parskip}{2pt}

\begin{document}

% ============================== PORTADA =====================================
\begin{titlepage}
\centering
\vspace*{1.2cm}
{\large UNIVERSIDAD TECNICA ESTATAL DE QUEVEDO}\\[0.25cm]
{\normalsize Facultad de Ciencias de la Computacion}\\[0.15cm]
{\normalsize Carrera de Ingenieria de Software}\\[1.6cm]

{\normalsize Aplicaciones Distribuidas (ISR-701) --- Unidad 5}\\[0.4cm]
{\large\bfseries FORM-TL-05}\\[0.3cm]
{\Large\bfseries Informe de Revision Tecnica Cruzada\\[0.2cm] del Codigo del PFC}\\[0.5cm]
{\normalsize Principios SOLID, patrones de diseno, separacion de capas,\\
manejo de excepciones distribuidas y observabilidad}\\[1.6cm]

\begin{tabularx}{0.92\textwidth}{@{}l X@{}}
\toprule
\multicolumn{2}{@{}l}{\textbf{Equipo revisor}}\\
\midrule
Sigla & FUVV\\
Integrante 1 & Freddy Vladimir Farinango Guandinango --- Bloques A y B\\
Integrante 2 & Isaias Abraham Urbina Romero --- Bloques C, D y E\\
\bottomrule
\end{tabularx}

\vspace{0.9cm}

\begin{tabularx}{0.92\textwidth}{@{}l X@{}}
\toprule
\multicolumn{2}{@{}l}{\textbf{PFC revisado (equipo ajeno)}}\\
\midrule
Sigla & ACC\\
Sistema & Soporte Tecnico ISP (gestion de tickets)\\
Repositorio & \url{\repoRevisado}\\
Commit auditado & \texttt{\commitCorto}\\
SHA completo & \texttt{\commitLargo}\\
Rama & \texttt{feature/entrega-3}\\
Fecha y hora & \fechaAuditoria\\
\bottomrule
\end{tabularx}

\vspace{1.2cm}
{\normalsize Repositorio del equipo revisor:}\\[0.15cm]
{\small\url{\repoRevisor}}

\vfill
{\small Periodo 2026--2027 \quad$\cdot$\quad Prof. PhD. Gleiston C. Guerrero-Ulloa}
\end{titlepage}

\tableofcontents
\newpage

% ============================ 1. OBJETO =====================================
\section{Objeto de la revision y congelacion del commit}

Este informe documenta la revision tecnica cruzada que el equipo FUVV realizo sobre el
Proyecto Fin de Curso del equipo ACC, \emph{Soporte Tecnico ISP}, conforme a la rotacion
fijada en la Seccion 2 de FORM-TL-05. Toda la revision se ejecuto sobre una unica
instantanea del repositorio, congelada antes de iniciar el recorrido de la lista de
verificacion. La salida literal de los comandos exigidos es la siguiente.

\begin{lstlisting}[caption={Comandos de verificacion del objeto revisado (Fase F1)},label={lst:git},numbers=none]
$ git log -1 --format="%H | %aI | %an"
479b6370a640c0d2edfad801e6937e9ede8c3c80 | 2026-07-27T21:51:36-05:00 | Jeremy Alvarez

$ git rev-parse --short HEAD
479b637

$ git ls-files | wc -l
1795

$ git rev-parse --abbrev-ref HEAD
feature/entrega-3
\end{lstlisting}

El repositorio acumula 34 commits y registra cuatro autores: Carlos Jose Carpio Mendoza,
Cristhian Daniel Pacheco Cardenas, Robinson Rodrigo Cando Moreno y Jeremy Alvarez. Todos
los hallazgos de este informe citan archivo, clase o funcion, numero de linea y el SHA
\texttt{\commitCorto}. Ninguna linea citada fue reconstruida de memoria: cada fragmento se
extrajo del arbol de trabajo correspondiente a ese commit.

\begin{tcolorbox}[colback=white,colframe=black!35,title=\textbf{Alcance y limitaciones declaradas}]
La revision es estatica. No se levanto el sistema completo (\texttt{start-all.ps1} requiere
un cluster CockroachDB de tres nodos, Kafka y MongoDB), por lo que los hallazgos relativos
a resiliencia (Bloque D) se sostienen en la lectura del codigo y de la configuracion, no en
una prueba de caos ejecutada por este equipo revisor. Alli donde el juicio depende de una
ejecucion no realizada, se indica explicitamente.
\end{tcolorbox}

% ======================== 2. ARQUITECTURA ===================================
\section{Mapa de arquitectura reconstruido}

El mapa siguiente se reconstruyo leyendo el arbol de fuentes, los ficheros
\texttt{application.yml}, los scripts de \ruta{db-cluster/scripts/} y las URL cableadas en
el frontend; no se tomo
de los diagramas incluidos en \texttt{docs/diagrams}, precisamente para poder contrastar
codigo contra documentacion.

\begin{table}[h]
\centering
\caption{Modulos desplegables identificados en el commit \texttt{\commitCorto}}
\label{tab:modulos}
\small
\begin{tabularx}{\textwidth}{@{}l l l X@{}}
\toprule
\textbf{Modulo} & \textbf{Tecnologia} & \textbf{Puerto} & \textbf{Responsabilidad observada en el codigo}\\
\midrule
\texttt{auth-service} & Spring Boot / Java & 8001 & Registro, login, refresh, logout, validacion de token y alta administrativa de usuarios.\\
\texttt{svc-principal} & Spring Boot / Java & 8002 & Ciclo de vida del ticket, autorizacion por rol y publicacion de eventos.\\
\texttt{report-service} & Spring Boot / Java & 8005 & Modelo de lectura CQRS reconstruido desde cuatro topicos y exportacion CSV.\\
\texttt{notification-service} & Node.js & --- & Consume tres topicos y persiste notificaciones simuladas en MongoDB.\\
\texttt{ai-service} & Python / FastAPI & --- & Clasificacion de tickets por reglas y publicacion de \texttt{ticket.classified}.\\
\texttt{frontend} & HTML/JS sin framework & --- & Cliente que invoca directamente los tres servicios HTTP.\\
\texttt{db-cluster} & CockroachDB & 26257 & Tres bases separadas: \texttt{auth\_db}, \texttt{ticket\_db}, \texttt{report\_db}.\\
\bottomrule
\end{tabularx}
\end{table}

\begin{figure}[h]
\centering
\begin{lstlisting}[numbers=none,basicstyle=\ttfamily\scriptsize]
   frontend/app.js  (3 URL cableadas: 8001, 8002, 8005)
        |               |                  |
        v               v                  v
  auth-service     svc-principal      report-service
    (8001)            (8002)              (8005)
        ^                |                  ^
        |  GET /validate |  publica         | consume 4 topicos
        +----------------+                  |
           (llamada sincrona por peticion)  |
                         |                  |
                         v                  |
                    +----------- KAFKA -----+-------------------+
                    | ticket.created / status-changed / assigned |
                    | ticket.classified                          |
                    +-----+--------------------------+-----------+
                          |                          |
                          v                          v
                    ai-service              notification-service
                    (clasifica)              (Mongo, simulado)

   Persistencia:  auth_db | ticket_db | report_db   (CockroachDB, 3 nodos)
                  NO hay base compartida entre servicios.
\end{lstlisting}
\caption{Mapa de dependencias reconstruido a partir del codigo. La unica arista sincrona
entre servicios es la validacion de token \texttt{svc-principal}$\rightarrow$\texttt{auth-service};
el resto de la integracion es asincrona por Kafka.}
\label{fig:arq}
\end{figure}

Dos observaciones estructurales se desprenden del mapa. Primera: no existe punto de entrada
unico; el navegador conoce y contacta tres servicios distintos, de modo que la funcion de
puerta de enlace esta ausente y sus responsabilidades transversales quedan duplicadas
dentro de cada servicio. Segunda: la separacion de datos por servicio si esta correctamente
implementada, lo que descarta de entrada el antipatron de base de datos compartida.

% ==================== 3. TABLA CONSOLIDADA ==================================
\section{Tabla consolidada de hallazgos}

\small
\begin{xltabular}{\textwidth}{@{}l l X X l@{}}
\caption{Hallazgos consolidados sobre el commit \texttt{\commitCorto} (Anexo B de FORM-TL-05)}\label{tab:hallazgos}\\
\toprule
\textbf{ID} & \textbf{Item} & \textbf{Hallazgo} & \textbf{Evidencia (ruta:linea)} & \textbf{Sev.}\\
\midrule
\endfirsthead
\toprule
\textbf{ID} & \textbf{Item} & \textbf{Hallazgo} & \textbf{Evidencia (ruta:linea)} & \textbf{Sev.}\\
\midrule
\endhead
\bottomrule
\endfoot
H01 & A1 & \texttt{TicketService} concentra negocio, autorizacion, persistencia, serializacion, mensajeria y metricas & \ruta{svc-principal/service/TicketService.java}:51-64, 99-128 & Mayor\\
H02 & A2 & La politica de roles se decide con cadenas de \texttt{if} replicadas en cuatro metodos & \ruta{TicketService.java}:74, 141-169, 226-253 & Mayor\\
H03 & A5 & La entidad de dominio \texttt{Ticket} importa Jakarta Persistence & \ruta{svc-principal/domain/Ticket.java}:3-9, 26-34 & Menor\\
H04 & A4 & \texttt{TicketRepository} hereda toda la superficie CRUD de \texttt{JpaRepository}, incluidos borrados que nadie usa & \ruta{svc-principal/repository/TicketRepository.java}:15-23 & Menor\\
H05 & B4/B5 & No existe puerta de enlace: el frontend cablea tres puertos y cada servicio duplica la validacion remota & \ruta{frontend/app.js}:5-7; \ruta{svc-principal/config/AuthGatewayFilter.java}:37-42 & Mayor\\
H06 & C2 & La capa de servicio importa DTO de la capa web & \ruta{TicketService.java}:12; \ruta{auth-service/service/AuthService.java}:14-18 & Menor\\
H07 & C5 & Secreto JWT y contrasena de administrador con valor por defecto versionado & \ruta{auth-service/resources/application.yml}:33, 42-43 & Mayor\\
H08 & C5 & URL de los tres servicios cableadas en el codigo del frontend & \ruta{frontend/app.js}:5-7; \ruta{frontend/auth/auth.js}:8 & Menor\\
H09 & D2/D4 & La llamada a \texttt{auth-service} se hace sin tiempo de espera explicito ni respuesta de degradacion & \ruta{svc-principal/config/AuthGatewayFilter.java}:40, 60-65, 76-78 & \sev{Critica}\\
H10 & D5 & El manejador generico devuelve \texttt{ex.getMessage()} al cliente & \ruta{svc-principal/web/GlobalExceptionHandler.java}:37-41; \ruta{AuthGatewayFilter.java}:77 & Mayor\\
H11 & E1/E2/E3 & Registro por consola en Node, sin identificador de correlacion ni formato estructurado & \ruta{notification-service/src/kafkaConsumer.js}:51, 56; \ruta{index.js}:23, 27 & Mayor\\
H12 & E4/E5 & Se registra el payload completo del evento y no existe auditoria de eventos criticos de autenticacion & \ruta{report-service/config/ReportEventListener.java}:82, 95, 105, 118; \ruta{AuthService.java}:40-54 & Mayor\\
\end{xltabular}

\normalsize
\noindent Los doce hallazgos se distribuyen en cuatro del Bloque A, uno del B, tres del C,
dos del D y dos del E, con una severidad Critica y siete Mayores. La columna de enlace al
issue correspondiente figura en la Tabla~\ref{tab:issues}.

% ======================= 4. ANALISIS POR BLOQUE =============================
\section{Analisis por bloque}

\begin{xltabular}{\textwidth}{@{}c c X@{}}
\caption{Lista de verificacion completa: 25 items marcados sobre el commit \texttt{\commitCorto}}\label{tab:checklist}\\
\toprule
\textbf{Item} & \textbf{Marca} & \textbf{Justificacion de la marca}\\
\midrule
\endfirsthead
\toprule
\textbf{Item} & \textbf{Marca} & \textbf{Justificacion de la marca}\\
\midrule
\endhead
\bottomrule
\endfoot
\multicolumn{3}{@{}l}{\textit{Bloque A --- Principios SOLID}}\\
A1 & N & \texttt{TicketService} acumula seis responsabilidades (H01).\\
A2 & N & Anadir un rol obliga a editar cuatro metodos existentes (H02).\\
A3 & NA & No hay jerarquias propias ni interfaces con multiples implementaciones: las unicas interfaces son los cuatro repositorios de Spring Data, cuya implementacion genera el contenedor. No existe ningun \texttt{UnsupportedOperationException} en codigo propio, de modo que no hay contrato de autor que pueda romperse.\\
A4 & P & Los repositorios exponen una interfaz mucho mas ancha que la usada por sus clientes (H04).\\
A5 & N & El modelo de dominio depende del ORM (H03).\\
\addlinespace
\multicolumn{3}{@{}l}{\textit{Bloque B --- Patrones de diseno}}\\
B1 & C & Se identifican Repository, CQRS y Saga por coreografia con sus roles asignados.\\
B2 & C & Cada patron responde a un problema real y esta justificado en ADR-0003 y ADR-0004.\\
B3 & P & \texttt{dispatcher.js} decide el canal con \texttt{switch} sobre el tipo de evento.\\
B4 & P & Base de datos por servicio correcta; puerta de enlace ausente (H05).\\
B5 & N & Cadena sincrona sin proteccion entre \texttt{svc-principal} y \texttt{auth-service} (H05, H09).\\
\addlinespace
\multicolumn{3}{@{}l}{\textit{Bloque C --- Separacion de capas}}\\
C1 & C & Los paquetes \texttt{web}, \texttt{service}, \texttt{domain}, \texttt{repository} y \texttt{config} existen como estructura real de paquetes.\\
C2 & N & La capa de servicio importa DTO de la capa web (H06); el dominio importa el ORM (H03).\\
C3 & C & \texttt{TicketResponse} y \texttt{SummaryResponse} median entre la entidad y la respuesta HTTP.\\
C4 & C & \texttt{TicketController} solo traduce y delega; la decision de rol vive en el servicio.\\
C5 & N & Secretos por defecto versionados (H07) y URL cableadas en el frontend (H08).\\
\addlinespace
\multicolumn{3}{@{}l}{\textit{Bloque D --- Excepciones distribuidas}}\\
D1 & P & No hay capturas silenciosas, pero \ruta{ReportEventListener.java}:131 registra el conflicto sin adjuntar la excepcion.\\
D2 & N & \texttt{RestClient.create()} sin tiempo de espera explicito (H09).\\
D3 & C & Reintento acotado a una sola vez sobre operacion repetible, con metrica asociada.\\
D4 & N & No hay degradacion: la caida de \texttt{auth-service} se traduce en 401 al usuario (H09).\\
D5 & N & Se filtra el mensaje interno de la excepcion al cliente (H10).\\
\addlinespace
\multicolumn{3}{@{}l}{\textit{Bloque E --- Logging y trazabilidad}}\\
E1 & P & Java y Python usan biblioteca de registro; Node imprime por consola (H11).\\
E2 & N & No se genera ni propaga identificador de correlacion en ningun punto del sistema.\\
E3 & N & Los mensajes son cadenas concatenadas, sin campos filtrables.\\
E4 & N & Se registra el payload completo del evento, con \texttt{clientId} y descripcion (H12).\\
E5 & N & \texttt{AuthService} no declara registrador ni componente de auditoria.\\
\end{xltabular}

\noindent Los cinco bloques se recorrieron completos: 25 items marcados, con la unica marca
NA justificada en la propia tabla. Los cinco principios SOLID fueron verificados. El resto de
esta seccion desarrolla la evidencia de cada hallazgo.
